\section{Выводы}
Выполнив восьмую лабораторную работу по курсу \enquote{Дискретный анализ}, я
познакомился с жадными алгоритмами и применил их к задаче о выборе самого
дешёвого набора пищевых добавок, по измерениям которых можно найти неизвестные коэффициенты.

Главное, что я понял: жадный алгоритм работает не везде, а только когда у задачи
есть подходящая структура. Здесь она появляется после сведения задачи к линейной алгебре.
Коэффициенты можно определить тогда и только тогда, когда ранг матрицы выбранных
добавок равен числу веществ $N$. Тогда искомый набор является базисом минимальной
стоимости, а множества линейно независимых строк образуют матроид, для
которого жадный выбор по возрастанию цены оптимален \cite{Kormen}. Без этого
наблюдения жадную сортировку по цене нельзя было бы обосновать, и мне пришлось бы перебирать подмножества.

На практике я научился проверять линейную независимость строки относительно
уже выбранных методом Гаусса, храня для каждого столбца индекс базисной строки.
Здесь важен порядок обработки: внутренние циклы идут от последнего столбца к текущему,
чтобы элемент, на который умножаются и делятся остальные, изменялся последним.
Ещё одна техническая деталь~--- сортировка подсчётом: она устойчива, работает
за $O(M + C)$ при ограниченных ценах и позволяет сохранить исходные номера добавок, нужные для вывода.

Слабое место решения~--- вычисления в числах с плавающей точкой: ноль определяется
сравнением с $\varepsilon = 10^{-10}$, поэтому при больших $N$ ошибки округления
могут исказить определение ранга. Для строгой гарантии пришлось бы использовать
вычисления в дробях или по модулю большого простого числа. Кроме того, тесты на
больших $N$ использовали случайные данные, поэтому я не могу утверждать, что на
специально подобранных тестах погрешности не проявятся.

\pagebreak